% Options for packages loaded elsewhere
\PassOptionsToPackage{unicode}{hyperref}
\PassOptionsToPackage{hyphens}{url}
%
\documentclass[
  brazilian,
]{article}
\usepackage{lmodern}
\usepackage{amssymb,amsmath}
\usepackage{ifxetex,ifluatex}
\ifnum 0\ifxetex 1\fi\ifluatex 1\fi=0 % if pdftex
  \usepackage[T1]{fontenc}
  \usepackage[utf8]{inputenc}
  \usepackage{textcomp} % provide euro and other symbols
\else % if luatex or xetex
  \usepackage{unicode-math}
  \defaultfontfeatures{Scale=MatchLowercase}
  \defaultfontfeatures[\rmfamily]{Ligatures=TeX,Scale=1}
\fi
% Use upquote if available, for straight quotes in verbatim environments
\IfFileExists{upquote.sty}{\usepackage{upquote}}{}
\IfFileExists{microtype.sty}{% use microtype if available
  \usepackage[]{microtype}
  \UseMicrotypeSet[protrusion]{basicmath} % disable protrusion for tt fonts
}{}
\makeatletter
\@ifundefined{KOMAClassName}{% if non-KOMA class
  \IfFileExists{parskip.sty}{%
    \usepackage{parskip}
  }{% else
    \setlength{\parindent}{0pt}
    \setlength{\parskip}{6pt plus 2pt minus 1pt}}
}{% if KOMA class
  \KOMAoptions{parskip=half}}
\makeatother
\usepackage{xcolor}
\IfFileExists{xurl.sty}{\usepackage{xurl}}{} % add URL line breaks if available
\IfFileExists{bookmark.sty}{\usepackage{bookmark}}{\usepackage{hyperref}}
\hypersetup{
  pdftitle={Parte III: Processamento de Linguagem Natural para Linguistas},
  pdfauthor={CiberExt 26-29 · FEELT38103 · Universidade Federal de Uberlândia},
  pdflang={pt-BR},
  hidelinks,
  pdfcreator={LaTeX via pandoc}}
\urlstyle{same} % disable monospaced font for URLs
\setlength{\emergencystretch}{3em} % prevent overfull lines
\providecommand{\tightlist}{%
  \setlength{\itemsep}{0pt}\setlength{\parskip}{0pt}}
\setcounter{secnumdepth}{-\maxdimen} % remove section numbering
\ifxetex
  % Load polyglossia as late as possible: uses bidi with RTL langages (e.g. Hebrew, Arabic)
  \usepackage{polyglossia}
  \setmainlanguage[]{brazilian}
\else
  \usepackage[shorthands=off,main=brazilian]{babel}
\fi

\title{Parte III: Processamento de Linguagem Natural para Linguistas}
\usepackage{etoolbox}
\makeatletter
\providecommand{\subtitle}[1]{% add subtitle to \maketitle
  \apptocmd{\@title}{\par {\large #1 \par}}{}{}
}
\makeatother
\subtitle{Unidade 3 - Exercícios}
\author{CiberExt 26-29 · FEELT38103 · Universidade Federal de
Uberlândia}
\date{Agosto de 2026}

\begin{document}
\maketitle

\hypertarget{exercuxedcios-encontrando-padruxf5es-linguuxedsticos-com-o-spacy}{%
\section{Exercícios --- Encontrando padrões linguísticos com o
spaCy}\label{exercuxedcios-encontrando-padruxf5es-linguuxedsticos-com-o-spacy}}

Os exercícios abaixo pressupõem que você tenha executado o código da
unidade e que as seguintes variáveis estejam disponíveis: \texttt{nlp},
\texttt{doc}, \texttt{matcher}, \texttt{morph\_matcher} e
\texttt{dep\_matcher}.

Para os exercícios de programação, use um modelo pequeno para o inglês
(\texttt{en\_core\_web\_sm}) e o arquivo \texttt{data/occupy.txt}
utilizado na unidade.

\hypertarget{parte-a-conceitos}{%
\subsection{Parte A --- Conceitos}\label{parte-a-conceitos}}

\textbf{Exercício 1.} Explique, com suas próprias palavras, a diferença
entre as classes \emph{Matcher}, \emph{DependencyMatcher} e
\emph{PhraseMatcher} do spaCy. Dê um exemplo de tarefa linguística
adequada a cada uma.

\textbf{Exercício 2.} No formato de padrões do spaCy, qual é o papel da
chave \texttt{OP}? Descreva o efeito de cada um dos quatro operadores
(\texttt{!}, \texttt{?}, \texttt{+}, \texttt{*}) e dê um exemplo de
padrão em que o operador \texttt{?} seria preferível ao \texttt{+}.

\textbf{Exercício 3.} Qual a diferença entre as chaves \texttt{POS} e
\texttt{TAG} em uma regra de padrão? Por que a etiqueta \texttt{NNP} não
aparece como valor de \texttt{POS}?

\textbf{Exercício 4.} Explique o funcionamento do atributo
\texttt{IS\_SUPERSET} aplicado à chave \texttt{MORPH}. Por que ele é
necessário quando queremos casar apenas \emph{alguns} traços
morfológicos?

\textbf{Exercício 5.} No \emph{DependencyMatcher}, explique a diferença
entre as chaves \texttt{LEFT\_ID} e \texttt{RIGHT\_ID}. Por que o
primeiro dicionário de um padrão não possui \texttt{LEFT\_ID}?

\textbf{Exercício 6.} Por que o \emph{DependencyMatcher} devolve tuplas
de índices em vez de objetos \emph{Span}, como faz o \emph{Matcher}?

\hypertarget{parte-b-programauxe7uxe3o}{%
\subsection{Parte B --- Programação}\label{parte-b-programauxe7uxe3o}}

\textbf{Exercício 7.} Crie um \emph{Matcher} que encontre sequências de
um adjetivo (\texttt{ADJ}) seguido de um substantivo (\texttt{NOUN}) no
objeto \texttt{doc}. Nomeie o padrão como \texttt{adj+noun} e imprima as
20 primeiras correspondências.

\textbf{Exercício 8.} Modifique o padrão do exercício anterior para
permitir \textbf{um ou mais} adjetivos antes do substantivo. Compare a
quantidade de correspondências obtidas com e sem o argumento
\texttt{greedy=\textquotesingle{}LONGEST\textquotesingle{}} e explique a
diferença.

\textbf{Exercício 9.} Escreva um padrão que case substantivos no plural,
usando a chave \texttt{MORPH} com \texttt{IS\_SUPERSET} e o traço
\texttt{Number=Plur}. Imprima as 15 primeiras correspondências junto com
os traços morfológicos de cada \emph{Token}.

\textbf{Exercício 10.} Usando o \emph{DependencyMatcher}, escreva um
padrão que encontre verbos (\texttt{VERB}) e seus objetos diretos
(\texttt{dobj}), sem incluir o sujeito. Nomeie o padrão como
\texttt{verb\_dobj} e imprima as 20 primeiras correspondências no
formato \texttt{verbo\ ...\ objeto}.

\textbf{Exercício 11.} Estenda o padrão do exercício anterior
acrescentando um terceiro elo à corrente: o \textbf{determinante}
(\texttt{det}) do objeto direto. Note que esse novo elo deve partir do
padrão \texttt{d\_object}, e não da âncora \texttt{verb}. Imprima as
correspondências no formato \texttt{verbo\ ...\ determinante\ objeto}.

\textbf{Exercício 12.} Escreva um laço que produza linhas de
concordância para as correspondências do exercício 7, exibindo 5
\emph{Tokens} antes e 5 \emph{Tokens} depois de cada correspondência,
com o trecho casado destacado em azul usando a classe \texttt{Printer}
da wasabi.

\hypertarget{parte-c-reflexuxe3o}{%
\subsection{Parte C --- Reflexão}\label{parte-c-reflexuxe3o}}

\textbf{Exercício 13.} As linhas de concordância da unidade às vezes
``pulam'' uma ou duas linhas na saída. Explique a causa e proponha uma
solução em Python que evite esse comportamento.

\textbf{Exercício 14.} Suponha que você queira estudar a construção
``verbo + preposição + sintagma nominal'' (por exemplo, \emph{depend on
the government}) em um corpus. Descreva qual \emph{Matcher} você usaria
e esboce a regra de padrão correspondente, justificando suas escolhas.

\end{document}
